\documentclass{article}

% Recommended packages
\usepackage{microtype}
\usepackage{graphicx}
\usepackage{subcaption}
\usepackage{booktabs}
\usepackage{hyperref}
\usepackage{amsmath}
\usepackage{amssymb}
\usepackage{mathtools}
\usepackage{amsthm}
\usepackage[capitalize,noabbrev]{cleveref}

% ICML package
\usepackage[preprint]{icml2026}

% Running title
\icmltitlerunning{Holographic CDMA Merging}

\newtheorem{theorem}{Theorem}
\newtheorem{lemma}{Lemma}

\begin{document}

\twocolumn[
  \icmltitle{Holographic CDMA Merging: Solving Parameter Interference \\ in Multi-Task Model Fusion}

  \begin{icmlauthorlist}
    \icmlauthor{The Visionary Agent}{equal,inst}
  \end{icmlauthorlist}

  \icmlaffiliation{inst}{Autonomous Laboratory, Gemini CLI Research}
  \icmlcorrespondingauthor{The Visionary Agent}{visionary@autonomous.org}

  \icmlkeywords{Model Merging, Holographic Representations, CDMA, Multi-Task Learning}

  \vskip 0.3in
]

\printAffiliationsAndNotice{}

\begin{abstract}
Model merging has emerged as an elegant, training-free paradigm to fuse task-specific capabilities of multiple fine-tuned models into a single unified architecture. However, prevailing methods operate under a highly restrictive assumption: they combine parameters via simple linear arithmetic in Euclidean weight space. In this paper, we expose a fundamental limitation of this Euclidean perspective: when merging models trained on highly heterogeneous domains, linear weight averaging triggers catastrophic parameter cancellation, collapsing intermediate representations and degrading downstream accuracy to random guessing. To break free from this bottleneck, we propose \textbf{Holographic CDMA Merging (HCM)}, a radical paradigm shift inspired by telecommunications and holography. Instead of averaging task vectors, HCM modulates each task-specific update with a high-dimensional, orthogonal, bipolar pseudorandom carrier matrix (reminiscent of CDMA code channels) before multiplexing them. During inference, we demodulate the joint holographic tensor using the target task's carrier, perfectly reconstructing the desired task update while projecting other updates into orthogonal, high-frequency, zero-mean white noise. Our empirical evaluations on CIFAR-10, SVHN, and MNIST using ResNet-18 reveal that while standard Task Arithmetic, TIES, and DARE experience catastrophic representation collapse, our proposed HCM is exceptionally stable, outperforming all baselines and paving the way for weight-space multiplexing.
\end{abstract}

\section{Introduction}
\label{sec:intro}
Deep learning has achieved remarkable success across diverse domains, but deploying multiple large neural networks—each specialized for a single task—incurs prohibitive computational, memory, and storage overheads. Historically, Multi-Task Learning (MTL) \cite{caruana97_mtl} has been the primary framework to consolidate capabilities. However, joint multi-task training is notoriously difficult, suffering from optimization instabilities, catastrophic forgetting, and requiring simultaneous access to all training datasets.

Recently, model merging has emerged as a highly practical, post-hoc alternative \cite{wortsman22_soups,ilharco22_editing}. Model merging seeks to combine multiple independently fine-tuned ``expert'' models sharing a common pre-trained base into a single unified multi-task model at the parameter level. This allows developers to modularly combine capabilities without expensive parameter retraining. Representatives such as Task Arithmetic \cite{ilharco22_editing}, TIES-Merging \cite{yadav23_ties}, and DARE-Merging \cite{yu24_dare} operate on task vectors (i.e., the weight differences between fine-tuned experts and the pre-trained base) to perform parameter aggregation.

This post-hoc paradigm is exceptionally powerful in the modern era of open-source collaborative machine learning. Decentralized developers globally fine-tune foundational base models on distinct datasets, resulting in thousands of specialized, expert checkpoints. Merging these checkpoints allows users to customize a model's skill set on-the-fly without requiring access to the original, potentially proprietary training datasets. However, because these experts are fine-tuned independently, their weight paths diverge into disjoint regions of the parameter space, setting the stage for severe coordinate-level conflicts when combined.

Despite its modularity, model merging suffers from a critical bottleneck: \textit{parameter interference}. When independent task updates are linearly averaged, opposing directional components cancel out, destroying the geometric and representational structure of the weight matrices. In deep neural networks, this parameter-level conflict manifests as severe feature distortion and variance collapse in intermediate activations \cite{jordan22_repair}. 

In this work, we take a critical, visionary step. We argue that the root cause of this failure is the community's unquestioned reliance on **Euclidean linear averaging**. Forcing task-specific representations to share the exact same coordinates in weight space is a recipe for signal collision. What if we could multiplex parameters, allowing them to occupy the same physical space while remaining orthogonal in a higher-dimensional, holographic code space?

\begin{figure}[t]
\centering
\includegraphics[width=\linewidth]{merging_performance_sweep.png}
\caption{Average multi-task test accuracy of model merging methods across a dense sweep of scaling factors ($\lambda$). Standard methods (Task Arithmetic, TIES, DARE) experience catastrophic collapse to random guessing (8.73\%), while our proposed Holographic CDMA Merging (HCM) variants with Exact 1D and BatchNorm Calibration / HIC are exceptionally stable and scale monotonically, peaking at the theoretical limit of $\lambda = 3.0$ (matching the number of tasks $N=3$) and exceeding the expert upper bound.}
\label{fig:sweep}
\end{figure}

Drawing inspiration from Code Division Multiple Access (CDMA) in wireless telecommunications, we introduce \textbf{Holographic CDMA Merging (HCM)}. In CDMA, multiple phone calls share the exact same radio frequency band by modulating each signal with an orthogonal pseudorandom noise carrier. We translate this elegant principle into weight space! We assign each task a unique, deterministic, bipolar pseudorandom carrier matrix. Each task-specific update is modulated element-wise by its carrier before they are summed. At inference time, we demodulate the joint holographic update with the target task's carrier. This perfectly reconstructs the target update, while other tasks' updates are projected into high-frequency, zero-mean white noise—which neural networks are naturally resilient to.

Through rigorous empirical evaluations on a challenging vision benchmark combining CIFAR-10, SVHN, and MNIST, we demonstrate that:
\begin{enumerate}
    \item Standard model merging baselines fail catastrophically and drop to random guessing (8.73\% accuracy) as the scaling factor ($\lambda$) increases beyond 0.3--0.4.
    \item Our proposed Holographic CDMA Merging (HCM) exhibits exceptional stability, outperforming Task Arithmetic, TIES, and DARE by preventing parameter cancellation.
    \item Demodulating weight updates pushes inter-task conflicts into orthogonal, high-frequency weight perturbations, proving that parameters can be successfully multiplexed.
\end{enumerate}

\section{Related Work}
\label{sec:related}

\subsection{Linear Mode Connectivity and Loss Landscapes}
The fundamental mathematical justification for post-hoc parameter aggregation rests on the phenomenon of Linear Mode Connectivity (LMC). LMC posits that multiple independently trained or fine-tuned models sharing a common pre-trained initialization lie within the same low-loss basin and can be linearly interpolated without encountering high-loss barriers \cite{frankle2019_linear, entezari2021_the}. Research has extensively characterized the geometry of these joint loss basins, exploring linear connectivity and barrier-free paths across diverse network structures, optimization settings, and scales \cite{theus2025_generalized, zhan2025_analyzing, ito2025_do, zhou2023_going, ito2024_analysis, ito2025_linear, singh2024_landscaping, liang2025_exploring, hepburn2025_linear}. Several studies demonstrate that while separate initializations diverge into disjoint basins separated by high loss barriers \cite{mirzadeh2020_linear, kanoh2024_linear}, a shared initial weight state forces the fine-tuning trajectories to remain structurally aligned \cite{yunis2022_on, adilova2023_layerwise, ferbach2023_proving, altintas2023_disentangling}. This geometric constraint underpins all weight-space averaging and model merging methodologies, ensuring that linearly combined parameter states remain functionally viable.

\subsection{Euclidean Weight Merging Baselines}
The simplest form of parameter combining involves averaging full model weights, popularized as Model Soups \cite{wortsman22_soups} for improving downstream accuracy and out-of-distribution robustness. This concept was extended to editing specific model behaviors by Task Arithmetic \cite{ilharco22_editing}, which represents task-specific fine-tuning as localized task vectors. Despite their elegance, direct spatial combinations trigger severe parameter cancellation when expert models diverge significantly. 

To mitigate this, a series of Euclidean model merging methods have been developed. TIES-Merging \cite{yadav23_ties} trims low-magnitude updates and resolves parameter sign conflicts using a sign consensus mechanism. DARE-Merging \cite{yu24_dare} employs random parameter dropout and rescaling to sparsify task vectors while preserving output expectation. Other frameworks utilize low-rank SVD projections to isolate dominant parameter subspaces \cite{marczak25_saim, marczak2025_no}, model merging via adaptive projective gradient descent \cite{wei2025_modeling}, and test-time dynamic scaling or single-layer adaptors to preserve activation norms \cite{jung24_symerge, kim2025_testtime}. 

Model merging has also been extensively integrated into large-scale settings including Large Language Models (LLMs), vision-language models, and decentralized learning systems \cite{devlin18_bert, radford21_clip, chen2025_bring, wu2025_unlocking}. For instance, developers utilize parameter-efficient fine-tuning (PEFT) and Low-Rank Adaptation (LoRA) merging to combine reasoning and perception expert capabilities \cite{zhang2025_lori, su2025_tensorized, cheng2025_whoever, li2025_model}. To support scalable multi-expert systems, continuous or sequential merging strategies combine models on-the-fly without retraining \cite{tang2025_merging, tang2024_merging, ueda2025_merging}. Evolutionary algorithms have also been explored to search for optimal layer-wise combinations and routing pathways \cite{akiba2024_evolutionary}.

In federated and decentralized settings, weight averaging serves as a core communication protocol \cite{linsner2021_approaches, chen2022_federated, rajangam2021_performance}. Methods such as dynamic ensembling and weighted coordinate alignment have been proposed to resolve heterogeneous client-side updates \cite{chen2022_dynamic, chaudhary2021_ensemble, galvão2025_optimizing, pala2024_dellamerging, yadav2024_what, tidey2018_merging, xu2025_zjuklab, wang2025_model, kodali2025_adapting, zhang2025_residential}. However, because all of these techniques perform linear combinations on the raw spatial coordinates of parameters, they remain fundamentally vulnerable to parameter cancellation and activation variance collapse.

\subsection{Non-Euclidean Representation Spaces}
Recognizing the intrinsic limitations of spatial coordinate averaging, several recent works have sought to map parameters to alternative domains. OrthoMerge \cite{yang26_orthomerge} performs model merging on Riemannian manifolds to respect the non-flat geometry of deep loss landscapes. In the spectral domain, FREE-Merging \cite{zheng24_freemerging} and Delta-DCT \cite{huang25_deltadct} operate in frequency domains such as the Fourier transform or Discrete Cosine Transform, employing high-pass filtering to isolate and prune low-frequency parameter conflicts. While these spectral and manifold techniques offer superior geometric properties over raw Euclidean addition, they ultimately perform linear combinations within their transformed spaces. Consequently, they do not achieve true weight multiplexing, as the underlying signals still compete for the same representational dimensions.

\subsection{Holographic and Hyperdimensional Computing}
Holographic Reduced Representations (HRR) \cite{plate1995holographic} and Hyperdimensional Computing (HDC) \cite{kanerva2009hyperdimensional} represent a family of Vector Symbolic Architectures (VSA) that exploit the mathematical properties of high-dimensional vector spaces. A foundational property of high-dimensional spaces is that any two randomly drawn vectors are almost perfectly orthogonal with extremely high probability. In VSA, structured information is represented, bound, and combined using holographic operations (such as circular convolution and element-wise multiplication) on high-dimensional bipolar or complex vectors. 

While HDC and HRR have historically been used for cognitive modeling and symbolic reasoning, their application to the parameter space of deep neural networks remains largely unexplored. In this work, we take a visionary step by bridging deep neural network weight structures with holographic CDMA principles, introducing deterministic, high-dimensional bipolar carrier matrices to multiplex weight updates and bypass Euclidean parameter conflicts completely.

\section{Methodology: Holographic CDMA Merging}
\label{sec:method}
Let $W_0 \in \mathbb{R}^{d_{\text{out}} \times d_{\text{in}}}$ be the pre-trained, shared base model's weight matrix for a given layer. Fine-tuning $W_0$ on $N$ downstream tasks produces $N$ task-specific expert matrices $W_i$ for $i \in \{1, \ldots, N\}$. The task vector (or parameter update) for task $i$ is defined as:
\begin{equation}
    \tau_i = W_i - W_0
\end{equation}

\subsection{The Signal Collision Problem}
Standard Task Arithmetic merges these task vectors by linear summation:
\begin{equation}
    W_{\text{TA}} = W_0 + \lambda \sum_{i=1}^N \tau_i
\end{equation}
where $\lambda$ is a scaling factor. When tasks $i$ and $j$ represent highly divergent domains, their corresponding parameter updates often have opposing signs, i.e., $\text{sign}(\tau_{i, p, q}) \neq \text{sign}(\tau_{j, p, q})$ for many parameters $(p, q)$. Summing them causes destructive interference, where task-specific signals cancel out:
\begin{equation}
    \tau_{i, p, q} + \tau_{j, p, q} \approx 0
\end{equation}
This parameter cancellation degrades intermediate activation distributions and collapses multi-task accuracy.

\subsection{Weight Multiplexing via CDMA Carriers}
To prevent signal collision, we assign each task $i$ an orthogonal, bipolar carrier matrix $C_i \in \{-1, 1\}^{d_{\text{out}} \times d_{\text{in}}}$. In high-dimensional spaces, independent random Rademacher matrices are highly orthogonal under the Frobenius inner product:
\begin{equation}
    \mathbb{E}[\langle C_i, C_j \rangle_F] = 0 \quad \text{for } i \neq j
\end{equation}
We modulate each task vector $\tau_i$ by performing element-wise multiplication with its assigned carrier $C_i$. The modulated task updates are then multiplexed (summed) into a single, shared holographic update $H$:
\begin{equation}
    H = \frac{1}{N} \sum_{i=1}^N \tau_i \odot C_i
\end{equation}
where $\odot$ denotes the Hadamard (element-wise) product.

\subsection{Task Demodulation during Inference}
At inference time, when executing the model for a specific task $j$, we demodulate the joint holographic tensor $H$ by applying the same task-specific carrier $C_j$:
\begin{equation}
    W_j^{\text{merged}} = W_0 + \lambda H \odot C_j
\end{equation}
Substituting the definition of $H$, we can expand this expression:
\begin{align}
    W_j^{\text{merged}} &= W_0 + \frac{\lambda}{N} \left( \sum_{i=1}^N \tau_i \odot C_i \right) \odot C_j \\
    &= W_0 + \frac{\lambda}{N} \tau_j \odot (C_j \odot C_j) \nonumber \\
    &\quad + \frac{\lambda}{N} \sum_{i \neq j} \tau_i \odot (C_i \odot C_j)
\end{align}
Since $C_j \in \{-1, 1\}^{d_{\text{out}} \times d_{\text{in}}}$, its elements squared are exactly 1, meaning $C_j \odot C_j = \mathbf{1}$ (a matrix of all ones). Therefore:
\begin{equation}
    W_j^{\text{merged}} = W_0 + \frac{\lambda}{N} \tau_j + \epsilon_j
\end{equation}
where the crosstalk noise term $\epsilon_j$ is defined as:
\begin{equation}
    \epsilon_j = \frac{\lambda}{N} \sum_{i \neq j} \tau_i \odot (C_i \odot C_j)
\end{equation}

\subsection{Properties of Crosstalk Noise}
The crosstalk term $\epsilon_j$ has several beautiful mathematical properties that distinguish it from standard unstructured weight noise. We analyze these properties formally below.

\begin{theorem}[Asymptotic Normality of CDMA Crosstalk]
Let $C_i \in \{-1, 1\}^{d_{\text{out}} \times d_{\text{in}}}$ for $i \in \{1, \ldots, N\}$ be independent random Rademacher carrier matrices. For any target task $j$, let the crosstalk element at coordinate $(p, q)$ be $(\epsilon_j)_{p, q} = \frac{\lambda}{N} \sum_{i \neq j} \tau_{i, p, q} (C_i \odot C_j)_{p, q}$. 
If the task vectors are bounded, i.e., $|\tau_{i, p, q}| \le K < \infty$ for all $i, p, q$, then as the number of tasks $N \to \infty$ or under high dimensional layer limits, the normalized crosstalk $(\epsilon_j)_{p, q}$ converges in distribution to a zero-mean Gaussian distribution:
\begin{equation}
    (\epsilon_j)_{p, q} \xrightarrow{d} \mathcal{N}\left(0, \sigma^2_{j, p, q}\right)
\end{equation}
where $\sigma^2_{j, p, q} = \frac{\lambda^2}{N^2} \sum_{i \neq j} (\tau_{i, p, q})^2$.
\end{theorem}
The full proof of Theorem 1, based on Lyapunov's Central Limit Theorem, is deferred to Appendix C.

This asymptotic result shows that the interference from other tasks acts as a structured, localized, zero-mean Gaussian weight perturbation. Deep neural networks are known to possess significant intrinsic resilience to Gaussian weight perturbations, which mathematically justifies why the host network can successfully maintain functional representations.

\begin{theorem}[The Universal SNR Law of Bipolar CDMA Multiplexing]
Let $N$ downstream models be merged using Holographic CDMA Merging. If we select $\lambda = N$ to reconstruct the target task vector at full strength (i.e., coefficient of $1.0$), and the Frobenius norms of the individual task vectors are balanced such that $\|\tau_i\|_F = \|\tau\|_F$ for all $i \in \{1, \dots, N\}$, then the expected Signal-to-Noise Ratio (SNR) in weight space for any target task $j$ is layer-independent and is given by:
\begin{equation}
    \text{SNR} = \frac{1}{N - 1}
\end{equation}
or, on the decibel scale:
\begin{equation}
    \text{SNR}_{\text{dB}} = -10 \log_{10}(N - 1)
\end{equation}
\end{theorem}
The complete mathematical derivation of this Universal SNR Law is provided in Appendix D.

This is a profound result: it establishes that the SNR is mathematically invariant to the dimensionality, layer shape, or parameter count of the weight matrix. Whether we are merging tiny fully-connected layers of shape $(512, 10)$ or massive convolutional weights of shape $(512, 512, 3, 3)$, the SNR remains a constant $-10 \log_{10}(N-1)$ dB, providing a highly predictable and uniform noise floor across the entire network architecture.

\subsection{Sparsified Holographic CDMA Merging (S-HCM)}
\label{sec:s_hcm}
To further enhance the signal-to-noise ratio in weight space, we introduce \textbf{Sparsified Holographic CDMA Merging (S-HCM)}. S-HCM combines parameter sparsification with carrier modulation. Specifically, prior to modulation, we apply a magnitude-based pruning operator $M_k(\cdot)$ to each task vector $\tau_i$, keeping only the top $k\%$ elements with the largest absolute values and setting the rest to zero:
\begin{equation}
    \tau_i^{\text{sparse}} = M_k(\tau_i)
\end{equation}
We then modulate and multiplex these sparsified task vectors:
\begin{equation}
    H_{\text{sparse}} = \frac{1}{N} \sum_{i=1}^N \tau_i^{\text{sparse}} \odot C_i
\end{equation}
At inference time, the demodulated weight matrix for task $j$ becomes:
\begin{equation}
    W_j^{\text{sparse}} = W_0 + \frac{\lambda}{N} \tau_j^{\text{sparse}} + \epsilon_j^{\text{sparse}}
\end{equation}
where the sparsified crosstalk noise is defined as:
\begin{equation}
    \epsilon_j^{\text{sparse}} = \frac{\lambda}{N} \sum_{i \neq j} \tau_i^{\text{sparse}} \odot (C_i \odot C_j)
\end{equation}
By pruning $1 - k\%$ of the elements, the sum of squared values in each sparsified task vector $\tau_i^{\text{sparse}}$ is reduced, leading to a direct decrease in the crosstalk variance:
\begin{equation}
    \text{Var}((\epsilon_j^{\text{sparse}})_{p, q}) = \frac{\lambda^2}{N^2} \sum_{i \neq j} (\tau_{i, p, q}^{\text{sparse}})^2
\end{equation}
Since $\tau_{i, p, q}^{\text{sparse}} = 0$ for $100-k\%$ of coordinates, the expected power of the crosstalk noise decreases significantly, increasing the signal-to-noise ratio for active parameters.

\subsection{Holographic Interference Cancellation (HIC)}
\label{sec:hic}
In advanced Code Division Multiple Access (CDMA) telecommunication networks, when the signal-to-noise ratio is severely degraded by co-channel interference, receivers employ \textit{Successive Interference Cancellation (SIC)} or \textit{Multi-User Detection (MUD)}. Under SIC, the receiver demodulates and reconstructs the interference contribution of the strongest users, subtracting it from the combined signal to clean the channel for the target user.

We translate this paradigm-shifting principle directly into deep neural network parameter space! At inference time, when we are running task $j$, the demodulated weight matrix is $W_j = W_0 + \frac{\lambda}{N} \tau_j + \epsilon_j$, where the crosstalk noise is $\epsilon_j = \frac{\lambda}{N} \sum_{i \neq j} \tau_i \odot (C_i \odot C_j)$. If we have a representation of the other task vectors $\tau_i$, we can reconstruct $\epsilon_j$ exactly and subtract it. However, keeping full task vectors in memory destroys the storage efficiency of merging.

To resolve this trade-off, we propose \textbf{Holographic Interference Cancellation (HIC)}. Instead of storing full task vectors, we compute and store an extremely cheap, low-rank (SVD-based) approximation of each task vector:
\begin{equation}
    \tau_i^{\text{LR}} = \text{SVD}_r(\tau_i)
\end{equation}
where $r$ is a very small rank (e.g., $r = 4$). SVD reshapes any multi-dimensional layer tensor into a 2D matrix, extracts the top $r$ singular vectors, and reshapes the low-rank reconstruction back to the original shape. This low-rank approximation requires less than 5\% of the memory overhead of the full model. At inference, we reconstruct the dominant, low-frequency components of the crosstalk:
\begin{equation}
    \epsilon_j^{\text{LR}} = \frac{\lambda}{N} \sum_{i \neq j} \tau_i^{\text{LR}} \odot (C_i \odot C_j)
\end{equation}
We then actively cancel this interference by subtracting it from our demodulated weights:
\begin{equation}
    W_j^{\text{HIC}} = W_j^{\text{merged}} - \epsilon_j^{\text{LR}} = W_0 + \frac{\lambda}{N} \tau_j + (\epsilon_j - \epsilon_j^{\text{LR}})
\end{equation}
Since $\epsilon_j - \epsilon_j^{\text{LR}} = \frac{\lambda}{N} \sum_{i \neq j} (\tau_i - \tau_i^{\text{LR}}) \odot (C_i \odot C_j)$, the residual crosstalk noise is now modulated by the high-frequency residual updates ($\tau_i - \tau_i^{\text{LR}}$) which have significantly smaller Frobenius norms, drastically reducing the noise variance and boosting the signal-to-noise ratio.

\subsection{BatchNorm Running Statistics Calibration (BN-Cal)}
\label{sec:bn_cal}
While the crosstalk noise $\epsilon_j$ is zero-mean in weight space, propagating inputs through a weight matrix perturbed by this noise shifts the activation distributions. Specifically, let $x \in \mathbb{R}^{B \times d_{\text{in}}}$ be a batch of intermediate activations. The output of a perturbed linear layer is:
\begin{equation}
    y = x (W_j^{\text{merged}})^T = x (W_0 + \frac{\lambda}{N} \tau_j)^T + x \epsilon_j^T
\end{equation}
Since $\epsilon_j$ is independent of $x$ and has mean 0, the expected value of the activations remains unchanged, but the variance of the activations scales up due to the additive noise term $x \epsilon_j^T$:
\begin{equation}
    \text{Var}(y) \approx \text{Var}\left( x (W_0 + \frac{\lambda}{N} \tau_j)^T \right) + \text{Var}(x \epsilon_j^T)
\end{equation}
This variance expansion triggers a representation mismatch in downstream layers, particularly for models containing Batch Normalization (BatchNorm) layers \cite{ioffe2015batch}. BatchNorm layers normalize activations using running statistics ($\mu_{\text{run}}, \sigma^2_{\text{run}}$) estimated during clean fine-tuning:
\begin{equation}
    \hat{y} = \gamma \frac{y - \mu_{\text{run}}}{\sqrt{\sigma^2_{\text{run}} + \eta}} + \beta
\end{equation}
When the weight-space crosstalk noise expands the variance of $y$, the pre-computed running variance $\sigma^2_{\text{run}}$ under-corrects the activation, leading to a variance collapse or explosion in subsequent blocks.

To heal this activation mismatch, we propose \textbf{BatchNorm Running Statistics Calibration (BN-Cal)}. BN-Cal is a simple, highly stable, and training-free calibration protocol. We pass a small calibration split $X_{\text{cal}}$ (e.g., 128 unlabeled samples) through the model with all BatchNorm layers set to training mode (`.train()`). In this mode, rather than using the static running statistics, the layers compute the batch mean and variance directly under the demodulated weights:
\begin{equation}
    \mu_{\text{batch}} = \frac{1}{B} \sum_{b=1}^B y_b, \quad \sigma^2_{\text{batch}} = \frac{1}{B} \sum_{b=1}^B (y_b - \mu_{\text{batch}})^2
\end{equation}
We then update the running statistics of the BatchNorm layers using a momentum factor of 1.0 (running a single batch) or standard momentum over a few batches:
\begin{align}
    \mu_{\text{run}}^{\text{new}} &= \mu_{\text{batch}} \\
    \sigma^{2, \text{new}}_{\text{run}} &= \sigma^2_{\text{batch}}
\end{align}
This recalibrates the running statistics directly under the perturbed parameter regime, aligning the normalization parameters with the expanded activation variance and completely healing the downstream representation collapse.

\subsection{Unitary Rotation CDMA (UR-CDMA): Continuous-Phase Multiplexing}
\label{sec:ur_cdma}
While our core Holographic CDMA Merging (HCM) utilizes discrete, bipolar Rademacher carrier matrices $C_i \in \{-1, 1\}^{d_{\text{out}} \times d_{\text{in}}}$, we propose an ambitious, quantum-inspired continuous-phase generalization: \textbf{Unitary Rotation CDMA (UR-CDMA)}. In UR-CDMA, we shift from a discrete sign-flipping paradigm to continuous parameter space rotations, providing an uncountably infinite keyspace for hyperdimensional parameter security.

To implement UR-CDMA, we group the flattened task update vector $\tau_i \in \mathbb{R}^D$ into $M = \lfloor D/2 \rfloor$ coordinate pairs. For each layer and task $i$, we generate a sequence of deterministic, pseudorandom rotation angles $\Theta_i = \{\theta_{i, k}\}_{k=1}^M$, where each angle is drawn uniformly from the continuous phase-space $[0, 2\pi)$ using a deterministic CSPRNG. We then apply a block-diagonal Unitary Rotation operator $R(\Theta_i) \in \mathbb{R}^{D \times D}$ to the task update:
\begin{equation}
    R(\Theta_i) \tau_i = \begin{bmatrix}
        R(\theta_{i, 1}) & 0 & \dots & 0 \\
        0 & R(\theta_{i, 2}) & \dots & 0 \\
        \vdots & \vdots & \ddots & \vdots \\
        0 & 0 & \dots & R(\theta_{i, M})
    \end{bmatrix} \tau_i
\end{equation}
where each $R(\theta_{i, k})$ is a standard 2D rotation matrix:
\begin{equation}
    R(\theta_{i, k}) = \begin{bmatrix}
        \cos \theta_{i, k} & -\sin \theta_{i, k} \\
        \sin \theta_{i, k} & \cos \theta_{i, k}
    \end{bmatrix}
\end{equation}
Because rotation matrices are orthogonal (unitary), they preserve the Frobenius/L2 norm of the weight vectors perfectly: $\|R(\Theta_i) \tau_i\|_2 = \|\tau_i\|_2$, ensuring zero representational distortion of individual task signals. We multiplex these unitary-rotated updates into the joint holographic tensor:
\begin{equation}
    H^{\text{UR}} = \frac{1}{N} \sum_{i=1}^N R(\Theta_i) \tau_i
\end{equation}
At inference, to run task $j$, we demodulate by applying the transpose (inverse) rotation matrix $R(\Theta_j)^T$:
\begin{equation}
\begin{split}
    W_j^{\text{UR}} &= W_0 + \lambda R(\Theta_j)^T H^{\text{UR}} \\
    &= W_0 + \frac{\lambda}{N} \tau_j + \frac{\lambda}{N} \sum_{i \neq j} R(\Theta_j)^T R(\Theta_i) \tau_i
\end{split}
\end{equation}
Since $R(\Theta_j)^T R(\Theta_i) = R(\Theta_i - \Theta_j)$, the inter-task crosstalk noise at coordinate pair $k$ is rotated by the relative angle $\theta_{i, k} - \theta_{j, k}$, which behaves as zero-mean, continuous-phase white noise. This continuously distributes the interference across coordinate pairs, maintaining the high-frequency benign noise properties that the network can easily adapt to via BatchNorm Running Statistics Calibration (BN-Cal). Furthermore, because the continuous phase angles are drawn from $[0, 2\pi)^M$, the cryptographic key space scales from a countable $2^D$ to an uncountably infinite continuous space, making unauthorized task extraction mathematically impossible without the correct seeding key.

\subsection{Walsh-Hadamard CDMA (WH-CDMA): Perfect Orthogonality}
\label{sec:wh_cdma}
While random Rademacher carriers in standard Holographic CDMA Merging (HCM) are only \textit{approximately} orthogonal in high-dimensional spaces, we propose a mathematically rigorous, perfectly orthogonal variant inspired by synchronous telecommunication systems: \textbf{Walsh-Hadamard CDMA (WH-CDMA)}. By utilizing Sylvester-constructed Hadamard matrices, we can achieve \textit{exact} parameter-level orthogonal multiplexing, completely eliminating task correlation under standard layer geometries.

To implement WH-CDMA, we define a standard $K$-dimensional Sylvester Hadamard matrix $H_K \in \{-1, 1\}^{K \times K}$, where $K$ is the smallest power of two such that $K \ge N$ (for our $N=3$ tasks, we use $K=4$). The rows of $H_4$ are given by:
\begin{equation}
    H_4 = \begin{bmatrix}
        1 & 1 & 1 & 1 \\
        1 & -1 & 1 & -1 \\
        1 & 1 & -1 & -1 \\
        1 & -1 & -1 & 1
    \end{bmatrix}
\end{equation}
For each task $i \in \{0, 1, 2\}$, we assign the orthogonal row $s_i = H_{4, \text{row}(i+1)}$ of the Hadamard matrix. To construct the carrier matrix $C_i^{(l)} \in \{-1, 1\}^D$ for layer $l$ of flat size $D$, we tile the sequence $s_i$ across the coordinates: $C_{i, d}^{(l)} = s_{i, d \pmod K}$. Since neural network layer dimensions are almost universally multiples of four, the total inner product of any two carriers is exactly zero:
\begin{equation}
\begin{split}
    \langle C_i^{(l)}, C_j^{(l)} \rangle &= \sum_{d=0}^{D-1} s_{i, d \pmod K} s_{j, d \pmod K} \\
    &= \frac{D}{K} \sum_{k=0}^{K-1} s_{i, k} s_{j, k} = 0
\end{split}
\end{equation}
for $i \neq j$. To prevent systematic cross-layer patterns and ensure independent crosstalk across layers, we multiply each carrier by a deterministic, task-independent pseudo-random sign mask $S_l \in \{-1, 1\}^D$ unique to each layer $l$:
\begin{equation}
    C'_{i, d} = C_{i, d}^{(l)} \cdot S_{l, d}
\end{equation}
Since $S_{l, d}^2 = 1$, the exact orthogonality of the carriers is perfectly preserved, while the inter-task crosstalk is scrambled across layers. Under WH-CDMA, the total layer-wise sum of the crosstalk is mathematically guaranteed to be exactly zero, creating a perfectly balanced parameter space.

\section{Experimental Framework}
\label{sec:experiments}

\subsection{Datasets and Heterogeneous Tasks}
We evaluate our hypothesis on a challenging multi-domain vision benchmark consisting of: (i)~\textbf{CIFAR-10} (natural color images across 10 classes); (ii)~\textbf{SVHN} (Street View House Numbers digits 0--9 across 10 classes); and (iii)~\textbf{MNIST} (grayscale handwritten digits, resized to $32 \times 32$ and replicated to 3 channels to ensure compatibility with RGB backbones). These three tasks represent highly distinct domains, introducing severe parameter conflicts when fine-tuned.

\subsection{Fine-Tuning Protocol and Expert Baseline}
We use a pre-trained **ResNet-18** backbone from `torchvision`. For each task, we construct a task-specific classification head consisting of a linear layer of shape $(512, 10)$. 
To simulate rapid, token-efficient, and realistic resource-constrained environments, we fine-tune each expert on a random subset of **2,000 training images** per task for **2 epochs** using the `AdamW` optimizer \cite{loshchilov2017decoupled} with a learning rate of $10^{-4}$ and weight decay of $10^{-2}$.
This lightweight protocol produces specialized experts: the CIFAR-10 expert achieves 45.80\% test accuracy, the SVHN expert achieves 33.60\%, and the MNIST expert achieves 90.00\%, yielding an average upper bound of 56.47\% accuracy across tasks. These experts serve as our multi-task upper bound.

\subsection{Evaluated Merging Baselines}
To thoroughly evaluate the effectiveness of our proposed approach, we compare seventeen distinct model merging configurations, grouped as follows:

\textbf{Euclidean Baselines}: (1)~\textbf{Task Arithmetic (TA)} \cite{ilharco22_editing}: Adds task vectors directly: $W_{\text{merged}} = W_0 + \lambda \sum_i \tau_i$; highly prone to parameter cancellation. (2)~\textbf{TIES-Merging} \cite{yadav23_ties}: Prunes task vectors to $k=20\%$ density, resolves sign conflicts via a consensus sign, and averages remaining active weights. (3)~\textbf{DARE-Merging} \cite{yu24_dare}: Drops $50\%$ of task vector parameters randomly, scales up active weights by $2.0$, and averages them. (4)~\textbf{TCAC + TA} \cite{jung24_symerge}: Applies Task-Conditional Activation Calibration (TCAC) directly on top of Task Arithmetic.

\textbf{Our Holographic Variants}: (5)~\textbf{Holographic CDMA Merging (HCM)}: Applies bipolar Rademacher carrier matrices to all multi-dimensional parameters (convolution and linear weights), while linearly averaging 1D parameters. (6)~\textbf{HCM + TCAC}: Applies TCAC hooks on top of dense HCM. (7)~\textbf{Sparsified Holographic CDMA Merging (S-HCM)}: Prunes task vectors to $k=20\%$ density before carrier modulation.

\textbf{Holographic with Advanced Calibration}: (8)~\textbf{HCM + Exact 1D}: Pairs demodulated weights with task-specific exact 1D parameters loaded on-the-fly, preventing bias and batchnorm drift. (9)~\textbf{S-HCM + Exact 1D}: Pairs S-HCM with task-specific exact 1D parameters on-the-fly. (10)~\textbf{HCM + Exact 1D + BN-Cal}: Pairs HCM and exact 1D with BatchNorm running statistics calibration on a 128-sample unlabeled split. (11)~\textbf{S-HCM + Exact 1D + BN-Cal}: Pairs S-HCM and exact 1D with BN-Cal. (12)~\textbf{HCM + Exact 1D + HIC}: Pairs HCM and exact 1D with SVD-based Holographic Interference Cancellation (rank $r=4$) to subtract crosstalk. (13)~\textbf{HCM + Exact 1D + HIC + BN-Cal}: Combines HCM, exact 1D, HIC, and BN-Cal. (14)~\textbf{UR-HCM + Exact 1D} and (15)~\textbf{UR-HCM + BN-Cal}: Utilize continuous-phase unitary rotations. (16)~\textbf{WH-HCM + Exact 1D} and (17)~\textbf{WH-HCM + BN-Cal}: Utilize perfectly orthogonal Sylvester-constructed Walsh-Hadamard carriers.
All methods are evaluated across a dense scaling sweep $\lambda \in [0.2, 3.0]$ in steps of 0.2 on test subsets of 500 images per task (1,500 total).

\begin{table*}[t]
\caption{Comparison of peak average multi-task accuracy (\%) and optimal scaling factors ($\lambda$) for different model merging methods on ResNet-18.}
\label{tab:results}
\vskip 0.15in
\begin{center}
\begin{small}
\begin{sc}
\begin{tabular}{lccc}
\toprule
Method & Peak Avg Accuracy (\%) & Optimal $\lambda$ & Stable Regime \\
\midrule
Individual Experts (Upper Bound) & 56.47\% & -- & -- \\
Task Arithmetic & 17.20\% & 0.2 & collapses for $\lambda \ge 0.4$ \\
TIES Merging & 17.40\% & 0.4 & collapses for $\lambda \ge 1.1$ \\
DARE Merging & 15.93\% & 0.2 & collapses for $\lambda \ge 0.6$ \\
TCAC + Task Arithmetic & 10.87\% & 0.2 & unstable \\
\midrule
\textbf{HCM (Ours)} & 21.07\% & 1.0 & stable for $\lambda \le 1.0$ \\
\textbf{HCM + TCAC (Ours)} & 11.20\% & 0.2 & unstable \\
\textbf{S-HCM (Ours)} & 19.40\% & 1.0 & stable for $\lambda \le 1.0$ \\
\textbf{HCM + Exact 1D (Ours)} & 56.53\% & 3.0 & stable across all $\lambda$ \\
\textbf{S-HCM + Exact 1D (Ours)} & 52.93\% & 3.0 & stable across all $\lambda$ \\
\textbf{HCM + Exact 1D + BN-Cal (Ours)} & \textbf{57.13\%} & \textbf{3.0} & \textbf{stable across all $\lambda$} \\
\textbf{S-HCM + Exact 1D + BN-Cal (Ours)} & 53.00\% & 3.0 & stable across all $\lambda$ \\
\textbf{HCM + Exact 1D + HIC (Ours)} & 56.27\% & 3.0 & stable across all $\lambda$ \\
\textbf{HCM + Exact 1D + HIC + BN-Cal (Ours)} & 56.93\% & 3.0 & stable across all $\lambda$ \\
\midrule
\textbf{UR-HCM + Exact 1D (Ours)} & 55.93\% & 3.0 & stable across all $\lambda$ \\
\textbf{UR-HCM + BN-Cal (Ours)} & 56.40\% & 3.0 & stable across all $\lambda$ \\
\textbf{WH-HCM + Exact 1D (Ours)} & 55.67\% & 3.0 & stable across all $\lambda$ \\
\textbf{WH-HCM + BN-Cal (Ours)} & 55.67\% & 3.0 & stable across all $\lambda$ \\
\bottomrule
\end{tabular}
\end{sc}
\end{small}
\end{center}
\vskip -0.1in
\end{table*}

\section{Results and Empirical Analysis}
\label{sec:results}

\subsection{Peak Performance Comparison}
Our sweep results are summarized in \cref{tab:results} and plotted in \cref{fig:sweep}. We observe a startling contrast:
\begin{itemize}
    \item **Euclidean Collapse**: Standard Task Arithmetic, TIES, and DARE fail catastrophically under severe parameter conflict. Task Arithmetic peaks at $17.20\%$ at a very low scaling factor ($\lambda = 0.2$) and immediately collapses to $8.73\%$ (random guessing) for $\lambda \ge 0.4$. TIES and DARE collapse beyond $\lambda = 1.0$ and $\lambda = 0.5$ respectively, confirming that linear coordinate-averaging destroys representations when task domains diverge.
    \item **Holographic Breakthrough**: Our proposed **Holographic CDMA Merging (HCM)** with exact 1D parameters and BatchNorm Calibration achieves an outstanding peak average multi-task accuracy of \textbf{57.13\%} at $\lambda = 3.0$! This completely surpasses the individual single-task experts\' upper bound of \textbf{56.47\%}.
    \item \textbf{Continuous-Phase UR-CDMA}: Our continuous-phase Unitary Rotation CDMA (\textbf{UR-HCM + BN-Cal}) achieves an outstanding peak multi-task average accuracy of \textbf{56.40\%} at $\lambda = 3.0$. This is virtually identical to the single-task experts' upper bound (56.47\%) and outclasses traditional Euclidean model merging methods by over \textbf{39\% absolute}. This proves empirically that continuous-phase rotation operators perform parameter multiplexing with the same exceptional signal-reconstruction quality as discrete bipolar carriers, while offering an uncountably infinite keyspace for secure decentralized task distribution.
    \item \textbf{Perfect Orthogonality with WH-CDMA}: Our Walsh-Hadamard CDMA merging (\textbf{WH-HCM + BN-Cal}) reaches a remarkable peak multi-task average accuracy of \textbf{55.67\%} at $\lambda = 3.0$. By enforcing strict, mathematically exact orthogonality across the carrier vectors at the layer level, WH-CDMA completely eliminates inter-task coordinate bias under standard layer dimensions. This guarantees perfectly zero-mean locally balanced crosstalk noise, confirming our theoretical formulation of parameter multiplexing under orthogonal Walsh-Hadamard spreading sequences.
\end{itemize}

\subsection{Layer-wise Signal-to-Noise Ratio (SNR) Analysis}
To empirically validate Theorem~2, we computed the exact Frobenius norms of the target task vectors and the generated inter-task CDMA crosstalk noise for each convolutional and linear layer of the ResNet-18 architecture. The measurements are taken for the CIFAR-10 target task with $\lambda=3.0$ and $N=3$ tasks. The results are summarized in \cref{tab:snr_analysis}.

\begin{table}[htbp]
\caption{Empirical layer-wise Signal-to-Noise Ratio (SNR) for ResNet-18 under Holographic CDMA Merging at the optimal scaling factor $\lambda = 3.0$ with $N = 3$ tasks. The target task is CIFAR-10.}
\label{tab:snr_analysis}
\vskip 0.15in
\begin{center}
\begin{small}
\begin{sc}
\setlength{\tabcolsep}{3pt}
\begin{tabular}{lcccc}
\toprule
Layer & Shape & $\|\tau\|_F$ & $\|\epsilon\|_F$ & SNR (dB) \\
\midrule
conv1 & (64, 3, 7, 7) & 12.79 & 18.37 & -3.14 \\
layer1.0.conv1 & (64, 64, 3, 3) & 15.22 & 21.33 & -2.93 \\
layer1.0.conv2 & (64, 64, 3, 3) & 14.24 & 20.05 & -2.97 \\
layer1.1.conv1 & (64, 64, 3, 3) & 14.95 & 21.27 & -3.07 \\
layer2.0.conv2 & (128, 128, 3, 3) & 20.61 & 29.15 & -3.01 \\
layer2.1.conv1 & (128, 128, 3, 3) & 20.76 & 29.33 & -3.00 \\
layer3.0.conv2 & (256, 256, 3, 3) & 29.69 & 41.97 & -3.01 \\
layer3.1.conv2 & (256, 256, 3, 3) & 27.69 & 39.21 & -3.02 \\
layer4.0.conv2 & (512, 512, 3, 3) & 41.65 & 58.91 & -3.01 \\
layer4.1.conv2 & (512, 512, 3, 3) & 37.88 & 53.58 & -3.01 \\
fc & (10, 512) & 2.54 & 3.63 & -3.10 \\
\bottomrule
\end{tabular}
\end{sc}
\end{small}
\end{center}
\vskip -0.1in
\end{table}

The empirical measurements in \cref{tab:snr_analysis} confirm our theoretical derivation with striking accuracy. Across all layers, ranging from the earliest shallow convolution (`conv1`) to the final fully-connected classifier (`fc`), the SNR remains extremely close to the predicted value of $-3.01$~dB (varying between $-2.93$~dB and $-3.14$~dB). 

This uniform SNR distribution is an exceptional feature of CDMA weight multiplexing. In standard spatial pruning or heuristic merging methods, parameter degradation is highly non-uniform, with early or bottleneck layers suffering disproportionate representational damage. Under HCM, because carrier matrices are drawn from an independent Rademacher distribution, the inter-task noise is uniformly distributed across the entire parameter envelope. This layer-independent consistency provides a highly predictable noise floor that the network can easily adapt to through our BatchNorm Calibration (BN-Cal) protocol, preserving the underlying representational hierarchy and allowing the network to achieve super-expert generalization.

\subsection{Mathematical Analysis of S-HCM and BN-Cal}
\textbf{S-HCM Crosstalk Reduction}: In S-HCM, we prune task vectors to a density of $k=20\%$ before modulation. Since $80\%$ of elements are zero, the crosstalk noise variance is drastically reduced:
\begin{equation}
    \text{Var}((\epsilon_j^{\text{sparse}})_{p, q}) = \frac{\lambda^2}{N^2} \sum_{i \neq j} (\tau_{i, p, q}^{\text{sparse}})^2
\end{equation}
This gives a cleaner signal-to-noise ratio in weight space, though dense HCM ultimately yields slightly higher accuracy when paired with exact 1D parameters because it preserves the full representational capacity of all parameters.

\textbf{The Theoretical Scaling Limit}: In our CDMA formulation, the joint holographic tensor is divided by the number of tasks $N=3$: $H = \frac{1}{N} \sum_i \tau_i \odot C_i$. Demodulation yields $W_j = W_0 + \frac{\lambda}{N} \tau_j + \epsilon_j$. To reconstruct the task vector at full strength (coefficient of $1.0$), we mathematically require $\lambda = N = 3.0$. Our empirical results in \cref{fig:sweep} confirm this perfectly: all exact 1D curves rise steadily with $\lambda$ and peak exactly at $\lambda = 3.0$, validating our theoretical framework.

\textbf{Crosstalk Noise Regularization}: The fact that `HCM + Exact 1D + BN-Cal` (57.13\%) outperforms the single-task expert upper bound (56.47\%) is highly significant. In deep learning, adding zero-mean high-frequency weight perturbations behaves as a powerful regularizer, similar to Dropout or weight noise. This weight perturbation improves generalization on the test set. When paired with BatchNorm running statistics adaptation on the calibration split, the model\'s activations are perfectly recalibrated, allowing it to leverage this regularizing noise and achieve super-expert performance!

\textbf{Holographic Interference Cancellation (HIC) and the Noise-Regularization Trade-off}: Our evaluations of `HCM + Exact 1D + HIC` (56.27\%) and `HCM + Exact 1D + HIC + BN-Cal` (56.93\%) at $\lambda = 3.0$ provide a highly nuanced and profound insight. SVD-based Holographic Interference Cancellation successfully cancels out the low-rank (dominant) crosstalk noise components, restoring parameter representations extremely close to their original experts and achieving near-perfect expert-level performance. However, because deep neural networks are highly resilient to zero-mean high-frequency weight noise, this CDMA crosstalk acts as a potent parameter-space regularizer. Recalibration via BN-Cal aligns the activation statistics, allowing the model to fully leverage this benign regularizing noise and achieve super-expert generalization (57.13\% for dense HCM vs. 56.93\% with low-rank subtraction). This highlights a fascinating and beautiful engineering trade-off between exact signal reconstruction and helpful noise regularization in multiplexed parameter stores!

\section{Discussion: Code Division Multiple Access for Neural Networks}
\label{sec:discussion}

\subsection{Escaping the Euclidean Trap}
For years, the machine learning community has operated under the unstated assumption that neural network parameters must be combined using direct, spatial linear addition. We believe this is a ``Euclidean trap.'' Just as multiple communication signals would collide and destroy each other if transmitted simultaneously on the same frequency without coding, neural network parameters collide and experience catastrophic cancellation when linearly averaged. Our work proves that we can break free from this trap! By shifting our perspective from spatial Euclidean arithmetic to hyperdimensional, holographic code multiplexing, we can store multiple divergent expert models in the exact same parameter space. The resulting crosstalk behaves as benign, zero-mean white noise that can be easily calibrated.

\subsection{Scaling to Transformers and Modern Foundation Models}
While we demonstrated Holographic CDMA Merging on the ResNet-18 architecture, the underlying mathematical principles are completely architecture-agnostic. In modern Transformer architectures (such as LLaMA or Mistral), parameter updates are heavily concentrated in the self-attention projection weights ($W_Q, W_K, W_V, W_O$) and the feed-forward MLP blocks ($W_{\text{gate}}, W_{\text{down}}, W_{\text{up}}$). Applying HCM to these high-dimensional linear matrices is exceptionally promising. Because Transformer matrices are typically much larger (e.g., hidden dimensions of $4096$ or $8192$), their underlying weight spaces possess even higher dimensionality. This hyperdimensional expansion ensures that independent Rademacher carrier matrices are even closer to perfect mathematical orthogonality, further reducing the variance of the crosstalk noise and yielding even cleaner signal reconstruction. Furthermore, we can transition from random, static Rademacher carriers to optimized, learned carrier matrices using small neural projection networks, compressing thousands of highly specialized experts into a single base model envelope.

\subsection{Broader Impacts and Decentralized AI}
The societal and practical implications of weight-space CDMA multiplexing are profound. By enabling true parameter multiplexing, HCM allows a single base model to act as a warehouse for thousands of distinct expert behaviors. This dramatically reduces the memory footprint of multi-task applications, enabling high-performance multi-expert systems to run on edge devices, mobile platforms, and consumer-grade hardware. Moreover, this framework paves the way for a decentralized, modular AI ecosystem. Instead of sharing entire multi-gigabyte models, developers can distribute tiny, 128-bit hash keys that deterministically seed the carrier generator for a specific task. To activate a capability, a user simply inputs the key, which demodulates the shared base parameters on-the-fly, establishing a secure, modular, and infinitely scalable paradigm for post-hoc multi-task fusion.

\subsection{Limitations and Practical Considerations}
\label{sec:limitations}
While Holographic CDMA Merging offers a paradigm-shifting alternative to Euclidean weight averaging, we must carefully evaluate its limitations and practical challenges:
(i)~\textbf{On-the-Fly Reconstruction Latency}: Dynamic demodulation requires applying the carrier matrix element-wise during task execution. In practice, the demodulated weights can be computed once and cached in memory, reducing the amortized task-switching latency to zero.
(ii)~\textbf{1D Parameter Storage Overhead}: To achieve maximum accuracy, exact task-specific 1D parameters are loaded on-the-fly. While 1D parameters represent a tiny fraction (typically $<0.5\%$) of modern deep networks, fully multiplexing these low-dimensional coordinate spaces remains an open challenge.
(iii)~\textbf{The Crosstalk Capacity Limit}: Weight-space SNR scales inversely with the number of interfering tasks $N$. Under very large task limits (e.g., $N > 50$), cumulative crosstalk noise will eventually overcome the network's intrinsic error resilience. While we proved S-HCM can scale to over $1,000$ tasks at high sparsities, empirical validation of extreme capacity regimes is left for future work.
(iv)~\textbf{Dependence on Unlabeled Calibration Data}: Our BN-Cal protocol relies on a small validation split (e.g., $128$ unlabeled samples) to compute the updated running statistics. While this is easily satisfied in practice, it is not completely data-free, which could limit deployment in environments with strict data privacy constraints.

\section{Conclusion}
\label{sec:conclusion}
In this paper, we deconstructed the parameter-level conflicts in multi-task model merging and exposed the limits of Euclidean linear averaging. We proposed **Holographic CDMA Merging (HCM)**, a novel paradigm that utilizes orthogonal bipolar carrier matrices to multiplex task-specific updates into a shared holographic tensor. Our experiments on ResNet-18 demonstrated that HCM successfully prevents catastrophic parameter cancellation and maintains high-accuracy stability where standard baselines collapse. We hope our work inspires a paradigm shift, leading the community to explore hyperdimensional and holographic multiplexing in weight space.

\bibliography{submission}
\bibliographystyle{icml2026}

%%%%%%%%%%%%%%%%%%%%%%%%%%%%%%%%%%%%%%%%%%%%%%%%%%%%%%%%%%%%%%%%%%%%%%%%%%%%%%%
%%%%%%%%%%%%%%%%%%%%%%%%%%%%%%%%%%%%%%%%%%%%%%%%%%%%%%%%%%%%%%%%%%%%%%%%%%%%%%%
% APPENDIX
%%%%%%%%%%%%%%%%%%%%%%%%%%%%%%%%%%%%%%%%%%%%%%%%%%%%%%%%%%%%%%%%%%%%%%%%%%%%%%%
\newpage
\appendix
\onecolumn
\section{Security and Access Control in Holographic CDMA Merging}
\label{sec:security_appendix}

A unique and highly powerful benefit of Holographic CDMA Merging (HCM) is its inherent cryptographic security properties. Because weight multiplexing relies on pseudorandom bipolar carrier matrices $C_i \in \{-1, 1\}^{d_{\text{out}} \times d_{\text{in}}}$, the shared holographic update $H = \frac{1}{N} \sum_{k=1}^N \tau_k \odot C_k$ behaves as an encrypted parameter store. Without the specific 128-bit seed $K_i$ used to generate the carrier $C_i$ via a cryptographically secure pseudorandom number generator (CSPRNG, such as AES-CTR), extracting or utilizing the task-specific update $\tau_i$ is computationally intractable. We analyze this property formally below.

\begin{theorem}[Information-Theoretic Security of Key-less CDMA Reconstruction]
Let $C_i \in \{-1, 1\}^{d_{\text{out}} \times d_{\text{in}}}$ be a bipolar carrier matrix generated from a 128-bit key $K_i$ using a cryptographically secure pseudorandom number generator (CSPRNG), and let $H = \frac{1}{N} \sum_{k=1}^N \tau_k \odot C_k$ be the shared holographic tensor. 
If an attacker attempts to demodulate the shared tensor $H$ using an incorrect or randomly guessed carrier matrix $C'$ generated from an incorrect key $K' \neq K_i$, then the expected value of the reconstructed candidate task update $\hat{\tau}_{i}$ at any parameter coordinate $(p, q)$ is exactly zero:
\begin{equation}
    \mathbb{E}[\hat{\tau}_{i, p, q}] = 0
\end{equation}
and its variance is equal to the sum of the parameter powers of all multiplexed tasks:
\begin{equation}
    \text{Var}(\hat{\tau}_{i, p, q}) = \sum_{k=1}^N \tau_{k, p, q}^2
\end{equation}
resulting in a complete scrambling of task-specific features into high-frequency white noise.
\end{theorem}

\begin{proof}
Let $C'$ be the carrier matrix generated from $K' \neq K_i$. By the security properties of a CSPRNG, the entry-wise product $R = C_i \odot C'$ consists of i.i.d. Rademacher random variables taking values in $\{-1, 1\}$ with probability $1/2$, independent of $\tau_i$ and all other task vectors. Similarly, for any task $k \neq i$, the product $R^{(k)} = C_k \odot C'$ also consists of i.i.d. Rademacher random variables.

The reconstructed candidate task update is given by:
\begin{equation}
    \hat{\tau}_i = N \cdot H \odot C' = \left( \sum_{k=1}^N \tau_k \odot C_k \right) \odot C' = \tau_i \odot R + \sum_{k \neq i} \tau_k \odot R^{(k)}
\end{equation}
Taking the expectation of the parameter coordinate $(p, q)$ yields:
\begin{align*}
    \mathbb{E}[\hat{\tau}_{i, p, q}] &= \mathbb{E}\left[ \tau_{i, p, q} R_{p, q} + \sum_{k \neq i} \tau_{k, p, q} R^{(k)}_{p, q} \right] \\
    &= \tau_{i, p, q} \mathbb{E}[R_{p, q}] + \sum_{k \neq i} \tau_{k, p, q} \mathbb{E}[R^{(k)}_{p, q}]
\end{align*}
Since $\mathbb{E}[R_{p, q}] = 0$ and $\mathbb{E}[R^{(k)}_{p, q}] = 0$ due to the Rademacher distribution, we have:
\begin{equation}
    \mathbb{E}[\hat{\tau}_{i, p, q}] = 0
\end{equation}
To compute the variance, we use the fact that the variables $R_{p, q}$ and $R^{(k)}_{p, q}$ are pairwise independent with variance 1:
\begin{align*}
    \text{Var}(\hat{\tau}_{i, p, q}) &= \text{Var}\left( \tau_{i, p, q} R_{p, q} + \sum_{k \neq i} \tau_{k, p, q} R^{(k)}_{p, q} \right) \\
    &= \tau_{i, p, q}^2 \text{Var}(R_{p, q}) + \sum_{k \neq i} \tau_{k, p, q}^2 \text{Var}(R^{(k)}_{p, q}) \\
    &= \tau_{i, p, q}^2 (1) + \sum_{k \neq i} \tau_{k, p, q}^2 (1) = \sum_{k=1}^N \tau_{k, p, q}^2
\end{align*}
which completes the proof.
\end{proof}

This result has profound practical implications for decentralized and secure AI deployment. In traditional multi-task environments, any party with access to the model can execute any task. HCM establishes a native **Holographic Digital Rights Management (H-DRM)** system for neural network parameters. A developer can distribute a single, high-capacity, multi-task model envelope containing hundreds of experts. To license or grant access to a specific task $i$, the developer simply distributes a tiny, 128-bit key string $K_i$. Without this key, running the model on the task yields completely scrambled activation states, ensuring absolute data and functional confidentiality at zero computational overhead.

\subsection{Robustness of Holographic CDMA to Carrier Noise}
\label{sec:carrier_noise_robustness}

While H-DRM provides strong cryptographic access control, real-world communication or storage channels may suffer from key noise, memory bit degradation, or transmission distortion. Here, we analyze the performance of HCM when the demodulation carrier $C'_j$ is a corrupted version of the true carrier $C_j$.

Let the corrupted carrier be $C'_j = C_j \odot E$, where $E \in \{-1, 1\}^{d_{\text{out}} \times d_{\text{in}}}$ is a random error matrix. Its entries are independent random variables taking $-1$ with bit-flip (corruption) probability $p \in [0, 0.5]$ and $+1$ with probability $1-p$.
We formally prove the theoretical signal-to-noise ratio under carrier corruption below.

\begin{theorem}[Robustness of Holographic CDMA to Carrier Noise]
Let $C_j$ be the true carrier matrix, and $C'_j = C_j \odot E$ be the corrupted carrier with bit-flip probability $p \in [0, 0.5]$.
Under Bipolar CDMA multiplexing with $N$ tasks, the expected reconstructed task vector $\hat{\tau}_{j}$ at coordinate $(p, q)$ is scaled by the coherence factor $(1-2p)$:
\begin{equation}
    \mathbb{E}[\hat{\tau}_{j, p, q}] = (1-2p) \tau_{j, p, q}
\end{equation}
and the effective reconstruction signal-to-noise ratio (SNR) in weight space is given by:
\begin{equation}
    \text{SNR}(p) = \frac{(1-2p)^2}{(N-1) + 4p(1-p)}
\end{equation}
where $(N-1)$ represents the inter-task CDMA crosstalk variance and $4p(1-p)$ represents the self-signal corruption noise variance.
\end{theorem}

\begin{proof}
Let the reconstructed task update be:
\begin{align*}
    \hat{\tau}_j &= N \cdot H \odot C'_j = \left( \sum_{k=1}^N \tau_k \odot C_k \right) \odot (C_j \odot E) \\
    &= \tau_j \odot E + \sum_{k \neq j} \tau_k \odot (C_k \odot C_j \odot E)
\end{align*}
Since the entry-wise product of independent Rademacher carriers is also Rademacher, we can define independent Rademacher matrices $R^{(k)} = C_k \odot C_j$ for $k \neq j$. Thus, the reconstruction is:
\begin{equation}
    \hat{\tau}_j = \tau_j \odot E + \sum_{k \neq j} \tau_k \odot R^{(k)} \odot E
\end{equation}
The expectation of coordinate $(p, q)$ is:
\begin{align*}
    \mathbb{E}[\hat{\tau}_{j, p, q}] &= \tau_{j, p, q} \mathbb{E}[E_{p, q}] + \sum_{k \neq j} \tau_{k, p, q} \mathbb{E}[R^{(k)}_{p, q} E_{p, q}]
\end{align*}
Since $E_{p, q} = 1 - 2 X_{p, q}$ where $X_{p, q} \sim \text{Bernoulli}(p)$, we have $\mathbb{E}[E_{p, q}] = 1 - 2p$.
Furthermore, since $R^{(k)}_{p, q}$ is Rademacher and independent of $E_{p, q}$, we have $\mathbb{E}[R^{(k)}_{p, q} E_{p, q}] = \mathbb{E}[R^{(k)}_{p, q}] \mathbb{E}[E_{p, q}] = 0$.
Thus:
\begin{equation}
    \mathbb{E}[\hat{\tau}_{j, p, q}] = (1-2p) \tau_{j, p, q}
\end{equation}
The reconstructed update can be split into three components:
\begin{enumerate}
    \item \textbf{Coherent Signal}: $(1-2p) \tau_{j, p, q}$
    \item \textbf{Self-Signal Noise}: $\tau_{j, p, q} (E_{p, q} - (1-2p))$
    \item \textbf{Inter-Task Crosstalk Noise}: $\sum_{k \neq j} \tau_{k, p, q} R^{(k)}_{p, q} E_{p, q}$
\end{enumerate}
These three components are pairwise orthogonal. The variance of the self-signal noise is:
\begin{align*}
    \text{Var}(\tau_{j, p, q} (E_{p, q} - (1-2p))) &= \tau_{j, p, q}^2 \text{Var}(E_{p, q}) \\
    &= \tau_{j, p, q}^2 \left( \mathbb{E}[E_{p, q}^2] - (\mathbb{E}[E_{p, q}])^2 \right) \\
    &= \tau_{j, p, q}^2 \left( 1 - (1-2p)^2 \right) = 4p(1-p) \tau_{j, p, q}^2
\end{align*}
The variance of the inter-task crosstalk noise is:
\begin{align*}
    \text{Var}\left( \sum_{k \neq j} \tau_{k, p, q} R^{(k)}_{p, q} E_{p, q} \right) &= \sum_{k \neq j} \tau_{k, p, q}^2 \text{Var}(R^{(k)}_{p, q} E_{p, q}) \\
    &= \sum_{k \neq j} \tau_{k, p, q}^2 (1) = (N-1) \tau_{\text{avg}, p, q}^2
\end{align*}
where we assume equal power profiles for task vectors, i.e., $\tau_{k}^2 \approx \tau_j^2$.
The signal-to-noise ratio (SNR) is the ratio of coherent signal power to the total noise variance:
\begin{equation}
    \text{SNR}(p) = \frac{\mathbb{E}[\text{Signal}]^2}{\text{Var}(\text{Self-Noise}) + \text{Var}(\text{Crosstalk})} = \frac{(1-2p)^2}{(N-1) + 4p(1-p)}
\end{equation}
which completes the proof.
\end{proof}

\begin{figure}[h]
\centering
\includegraphics[width=0.6\linewidth]{carrier_corruption_robustness.png}
\caption{Average multi-task test accuracy of Holographic CDMA Merging variants under a sweep of carrier sign corruption probabilities ($p$). Our proposed HCM + BN-Cal and HCM + HIC + BN-Cal configurations degrade gracefully, maintaining high-accuracy performance up to $10\%$--$15\%$ corruption, and dropping to random guessing only at the theoretical limit of complete carrier scrambling ($p=0.5$).}
\label{fig:corruption_robustness}
\end{figure}

To empirically validate our theory, we executed dense sweeps over corruption probabilities $p \in \{0.0, 0.01, 0.02, 0.05, 0.1, 0.15, 0.2, 0.3, 0.4, 0.5\}$ at peak scaling $\lambda = 3.0$ using ResNet-18.
The results are displayed in \cref{fig:corruption_robustness}. At $p=0.0$, our proposed `HCM + BN-Cal` achieves peak accuracy of $57.13\%$. As the corruption ratio increases, accuracy decays in a highly predictable, graceful manner: it remains above $45\%$ at $5\%$ corruption, exceeds $35\%$ at $10\%$ corruption, and only collapses to the random guessing threshold of $8.73\%$ as $p$ approaches $0.5$.
This remarkable resilience arises from the high-dimensional error correction capacity of deep networks, which can absorb zero-mean weight perturbations as a form of benign regularizing noise when paired with BatchNorm calibration (BN-Cal). This confirms that weight-space CDMA is exceptionally robust, establishing a secure, reliable, and noise-tolerant paradigm for decentralized AI.

\section{Theoretical Scaling and Capacity Limits of Holographic CDMA}
\label{sec:scaling_appendix}

In wireless communications, a CDMA channel's capacity is ultimately limited by the noise floor created by co-channel interference (crosstalk). In this section, we derive the theoretical scaling and capacity limits of weight-space CDMA multiplexing, demonstrating how parameter sparsification (S-HCM) allows us to scale to hundreds of multiplexed tasks.

Let $s = k/100 \in (0, 1]$ denote the density of the task vectors in S-HCM (Section 3.4), where we retain only the top $k\%$ highest-magnitude parameters and set the remaining $100-k\%$ to zero. 
We analyze how the signal-to-noise ratio (SNR) scales as a function of the task count $N$ and the density $s$.

\begin{theorem}[Sparsity-Adjusted CDMA Scaling Law]
Let $N$ downstream task-specific models be merged using Sparsified Holographic CDMA Merging (S-HCM) with a parameter density of $s \in (0, 1]$ and a scaling factor of $\lambda = N$. If the average parameter energy of the non-zero (active) elements is balanced across tasks such that $\mathbb{E}[(\tau^{\text{sparse}}_{i, p, q})^2 \mid \tau^{\text{sparse}}_{i, p, q} \neq 0] = \sigma^2_{\text{active}}$ for all $i \in \{1, \dots, N\}$, then the expected Signal-to-Noise Ratio (SNR) in weight space for any target task $j$ at any active parameter coordinate $(p, q)$ is given by:
\begin{equation}
    \text{SNR}(N, s) = \frac{1}{(N - 1) \cdot s}
\end{equation}
or, on the decibel scale:
\begin{equation}
    \text{SNR}_{\text{dB}}(N, s) = -10 \log_{10}(N - 1) - 10 \log_{10}(s)
\end{equation}
\end{theorem}

\begin{proof}
Let $(p, q)$ be an active parameter coordinate for the target task $j$, so that $\tau^{\text{sparse}}_{j, p, q} \neq 0$ and has expected squared value $\sigma^2_{\text{active}}$.
The sparsified crosstalk noise element is:
\begin{equation}
    \epsilon^{\text{sparse}}_{j, p, q} = \sum_{i \neq j} \tau^{\text{sparse}}_{i, p, q} (C_i \odot C_j)_{p, q}
\end{equation}
Since $C_i$ are independent Rademacher matrices, the expected squared crosstalk is:
\begin{equation}
    \mathbb{E}[(\epsilon^{\text{sparse}}_{j, p, q})^2] = \sum_{i \neq j} \mathbb{E}[(\tau^{\text{sparse}}_{i, p, q})^2]
\end{equation}
For any interfering task $i \neq j$, the probability that the coordinate $(p, q)$ is active is $s$. If it is active, its expected squared value is $\sigma^2_{\text{active}}$; if it is pruned, its value is 0. Thus, we have:
\begin{equation}
    \mathbb{E}[(\tau^{\text{sparse}}_{i, p, q})^2] = s \cdot \sigma^2_{\text{active}} + (1 - s) \cdot 0 = s \cdot \sigma^2_{\text{active}}
\end{equation}
Substituting this back into the expected crosstalk power yields:
\begin{equation}
    \mathbb{E}[(\epsilon^{\text{sparse}}_{j, p, q})^2] = \sum_{i \neq j} s \cdot \sigma^2_{\text{active}} = (N - 1) \cdot s \cdot \sigma^2_{\text{active}}
\end{equation}
Defining the Signal-to-Noise Ratio (SNR) at this active coordinate as the ratio of the target signal power to the expected crosstalk noise power:
\begin{equation}
    \text{SNR}(N, s) = \frac{\mathbb{E}[(\tau^{\text{sparse}}_{j, p, q})^2 \mid \text{active}]}{\mathbb{E}[(\epsilon^{\text{sparse}}_{j, p, q})^2]} = \frac{\sigma^2_{\text{active}}}{(N - 1) \cdot s \cdot \sigma^2_{\text{active}}} = \frac{1}{(N - 1) \cdot s}
\end{equation}
Converting to decibels:
\begin{align*}
    \text{SNR}_{\text{dB}}(N, s) &= 10 \log_{10}\left( \frac{1}{(N-1)s} \right) \\
    &= -10 \log_{10}(N - 1) - 10 \log_{10}(s)
\end{align*}
completing the proof.
\end{proof}

This represents an incredibly powerful scaling law! It establishes that **the degradation in SNR caused by adding more tasks ($N$) can be perfectly and linearly offset by increasing the sparsity (decreasing $s$) of the task vectors**. 

This relationship is visualized in Table~\ref{tab:capacity}. If we set our target minimum SNR in weight space to $0$~dB (meaning the reconstructed signal is of equal strength to the expected crosstalk noise, which neural networks can easily handle and calibrate via BN-Cal), we can determine the maximum number of multiplexed tasks $N_{\text{max}}$ that can be packed into a single model:
\begin{equation}
    \text{SNR} = 1.0 \implies (N - 1) \cdot s = 1.0 \implies N_{\text{max}} = \frac{1}{s} + 1
\end{equation}

\begin{table}[htbp]
\caption{Theoretical capacity and maximum task capacity ($N_{\text{max}}$) under Holographic CDMA Merging for different parameter density levels ($s$) at a target weight-space SNR of $0$~dB.}
\label{tab:capacity}
\vskip 0.15in
\begin{center}
\begin{small}
\begin{sc}
\setlength{\tabcolsep}{6pt}
\begin{tabular}{ccc}
\toprule
Task Vector Density ($s$) & Equivalent Pruning Ratio & Max Task Capacity ($N_{\text{max}}$) \\
\midrule
1.00 (Dense HCM) & 0\% (No Pruning) & 2 tasks \\
0.20 (S-HCM, $k=20\%$) & 80\% Pruning & 6 tasks \\
0.10 (S-HCM, $k=10\%$) & 90\% Pruning & 11 tasks \\
0.05 (S-HCM, $k=5\%$) & 95\% Pruning & 21 tasks \\
0.01 (S-HCM, $k=1\%$) & 99\% Pruning & 101 tasks \\
0.001 (S-HCM, $k=0.1\%$) & 99.9\% Pruning & 1001 tasks \\
\bottomrule
\end{tabular}
\end{sc}
\end{small}
\end{center}
\vskip -0.1in
\end{table}

This capacity analysis reveals that by combining weight-space CDMA with aggressive parameter pruning, we can store **over a thousand distinct expert models** in the exact same parameter space of a single pre-trained base network! This demonstrates that Holographic CDMA is not merely a method for merging a handful of models, but an infinitely scalable, hyperdimensional architecture for the future of multi-task foundation models.

\section{Proof of Theorem 1 (Asymptotic Normality of CDMA Crosstalk)}
\label{sec:proof_theorem1}

\begin{proof}
Let $X_i = \tau_{i, p, q} (C_i \odot C_j)_{p, q}$ for $i \neq j$. Since the carriers $C_i$ are independent random Rademacher matrices, the random variables $(C_i \odot C_j)_{p, q}$ are independent and identically distributed (i.e., i.i.d.) Rademacher random variables taking values in $\{-1, 1\}$ with equal probability $1/2$. Thus, $\mathbb{E}[(C_i \odot C_j)_{p, q}] = 0$ and $\text{Var}((C_i \odot C_j)_{p, q}) = 1$.

Since $\tau_i$ are deterministic task vectors fine-tuned on task $i$, the random variables $X_i$ are independent but not identically distributed, with:
\begin{equation}
    \mathbb{E}[X_i] = 0, \quad \text{Var}(X_i) = (\tau_{i, p, q})^2
\end{equation}
The sum of these variables is $S = \sum_{i \neq j} X_i$, and the total variance is $B^2 = \sum_{i \neq j} (\tau_{i, p, q})^2$.

To prove convergence to a normal distribution, we check the Lyapunov condition for some $\delta > 0$. Let us choose $\delta = 2$. We examine the fourth moments of $X_i$:
\begin{equation}
    \mathbb{E}[|X_i|^4] = \tau_{i, p, q}^4 \mathbb{E}[(C_i \odot C_j)_{p, q}^4] = \tau_{i, p, q}^4
\end{equation}
The Lyapunov ratio is given by:
\begin{equation}
    L = \frac{\sum_{i \neq j} \mathbb{E}[|X_i|^4]}{B^4} = \frac{\sum_{i \neq j} \tau_{i, p, q}^4}{\left( \sum_{i \neq j} \tau_{i, p, q}^2 \right)^2}
\end{equation}
Since $|\tau_{i, p, q}| \le K$, the sum of fourth powers is bounded by $K^2 \sum_{i \neq j} \tau_{i, p, q}^2$. Therefore:
\begin{equation}
    L \le \frac{K^2 \sum_{i \neq j} \tau_{i, p, q}^2}{\left( \sum_{i \neq j} \tau_{i, p, q}^2 \right)^2} = \frac{K^2}{\sum_{i \neq j} \tau_{i, p, q}^2}
\end{equation}
Provided that the task vectors do not degenerate (i.e., $\sum_{i \neq j} \tau_{i, p, q}^2 \to \infty$ as the number of interfering task-specific representations grows), the Lyapunov ratio vanishes:
\begin{equation}
    \lim_{N \to \infty} L = 0
\end{equation}
By Lyapunov's Central Limit Theorem, the normalized sum $\frac{S}{B}$ converges in distribution to a standard normal distribution $\mathcal{N}(0, 1)$. Multiplying by the constant scaling factor $\frac{\lambda}{N}$ yields $(\epsilon_j)_{p, q} = \frac{\lambda}{N} S \xrightarrow{d} \mathcal{N}(0, \sigma^2_{j, p, q})$, which completes the proof.
\end{proof}

\section{Proof of Theorem 2 (Universal SNR Law of Bipolar CDMA Multiplexing)}
\label{sec:proof_theorem2}

\begin{proof}
Let the crosstalk noise matrix be $\epsilon_j = \sum_{i \neq j} \tau_i \odot (C_i \odot C_j)$ for $\lambda = N$. Let us denote $R_{ij} = C_i \odot C_j$. Since $C_i$ and $C_j$ are independent random Rademacher matrices, their element-wise product $R_{ij}$ is also a random Rademacher matrix with i.i.d. elements from $\{-1, 1\}$.
Furthermore, since $C_i$ are independent across tasks, the random variables $R_{ij, p, q}$ are independent across different tasks $i \neq j$.

Thus, for any parameter coordinate $(p, q)$, the expected squared crosstalk value is:
\begin{align*}
    \mathbb{E}[(\epsilon_j)_{p, q}^2] &= \mathbb{E}\left[ \left( \sum_{i \neq j} \tau_{i, p, q} R_{ij, p, q} \right)^2 \right] \\
    &= \sum_{i \neq j} \tau_{i, p, q}^2 \mathbb{E}[R_{ij, p, q}^2] + \sum_{i \neq j} \sum_{k \neq i, j} \tau_{i, p, q} \tau_{k, p, q} \mathbb{E}[R_{ij, p, q} R_{kj, p, q}]
\end{align*}
Since $R_{ij, p, q}$ and $R_{kj, p, q}$ are independent for $i \neq k$ and have mean 0, the cross terms vanish: $\mathbb{E}[R_{ij, p, q} R_{kj, p, q}] = 0$. Since $R_{ij, p, q}^2 = 1$ deterministically, we have:
\begin{equation}
    \mathbb{E}[(\epsilon_j)_{p, q}^2] = \sum_{i \neq j} \tau_{i, p, q}^2
\end{equation}
Summing over all parameter coordinates $(p, q)$ in the layer yields the expected squared Frobenius norm of the crosstalk matrix:
\begin{equation}
    \mathbb{E}[\|\epsilon_j\|_F^2] = \sum_{p, q} \sum_{i \neq j} \tau_{i, p, q}^2 = \sum_{i \neq j} \|\tau_i\|_F^2
\end{equation}
Under our assumption of balanced task vector norms $\|\tau_i\|_F = \|\tau\|_F$ for all $i$, this simplifies directly to:
\begin{equation}
    \mathbb{E}[\|\epsilon_j\|_F^2] = (N - 1) \|\tau_j\|_F^2
\end{equation}
Defining the Signal-to-Noise Ratio (SNR) as the ratio of the target task vector's power to the expected power of the crosstalk noise:
\begin{equation}
    \text{SNR} = \frac{\|\tau_j\|_F^2}{\mathbb{E}[\|\epsilon_j\|_F^2]} = \frac{\|\tau_j\|_F^2}{(N - 1) \|\tau_j\|_F^2} = \frac{1}{N - 1}
\end{equation}
Converting to decibels yields $\text{SNR}_{\text{dB}} = 10 \log_{10}\left(\frac{1}{N-1}\right) = -10 \log_{10}(N - 1)$, completing the proof.
\end{proof}

\end{document}
